\chapter*{前言}
\addcontentsline{toc}{chapter}{前言}

\section*{本书的目标}

这本书不是一本传统的数字电路教材。

传统教材的写法是：数制→逻辑门→布尔代数→卡诺图→组合逻辑→时序逻辑→... \par

这本书的写法是：\par
\begin{center}
\fbox{%
\begin{minipage}{0.85\textwidth}
\textbf{以电路为主线，反向展开其所涉及的数字电路知识。}

电路 → 原理 → 状态变化 → 设计思想 → 设计模板 → VHDL实现 → 考试变式
\end{minipage}%
}
\end{center}

\section*{本书的结构}

本书分为十个部分：

\begin{description}
  \item[第一部分] 组合逻辑电路体系 —— 编码器、译码器、MUX、比较器、加法器
  \item[第二部分] D触发器体系 —— 从"为什么需要存储"开始，建立状态与记忆的概念
  \item[第三部分] 计数器体系 —— 用统一框架串起所有模值计数器
  \item[第四部分] 分频器体系 —— 从状态变化解释频率减半的物理过程
  \item[第五部分] 序列发生器体系 —— 多种实现方案比较
  \item[第六部分] 移位寄存器体系 —— 数据流动的逐拍分析
  \item[第七部分] 环形计数器体系 —— 从移位寄存器推导环形计数器
  \item[第八部分] 序列检测器体系 —— 状态机的核心应用
  \item[第九部分] VHDL专题 —— 每种VHDL写法综合成什么硬件
  \item[第十部分] 考试训练体系 —— 每类电路20+道变式题
\end{description}

\section*{如何使用本书}

\begin{enumerate}
  \item \textbf{如果你时间充裕}：从头到尾，按顺序阅读。
  \item \textbf{如果你时间紧张}：直接跳到对应电路的章节，每章是自包含的。
  \item \textbf{如果你只需要刷题}：直接看第十部分。
\end{enumerate}

\section*{关于"为什么"}

本书最重要的一个原则是：

\begin{center}
\fbox{%
\begin{minipage}{0.85\textwidth}
\large \textbf{不讲"是什么"，只讲"为什么"。}

每一个设计决策都有原因。

每一个公式都有推导过程。

每一个状态变化都有逐拍分析。
\end{minipage}%
}
\end{center}

\section*{关于代码}

\textbf{不要在理解电路之前看代码。}

如果你看到一个电路，第一反应是"它的VHDL怎么写"，那你就学反了。

正确的顺序是：

\begin{enumerate}
  \item 理解这个电路要解决什么问题
  \item 理解它的输入输出关系
  \item 理解它的内部工作原理
  \item 理解为什么这样设计
  \item 然后，VHDL只是这些理解的翻译
\end{enumerate}

如果你理解了设计逻辑，任何VHDL变式你都能自己写出来。

\section*{写在前面}

这套书可能和你在学校里见过的任何一本数字电路教材都不一样。

它不讲数制，不讲卡诺图，不从逻辑代数开始。(它默认你已经学过这些了，并且已经能够利用这些工具分析电路了。)

它从数字电路中典型的23个电路开始。

每一个电路，我们不讲废话，只讲三件事：

\begin{enumerate}
  \item \textbf{它为什么存在}（解决了什么问题）
  \item \textbf{它是怎么工作的}（原理和状态变化）
  \item \textbf{怎么用它解题}（考试变式和解题模板）
\end{enumerate}

祝学习顺利。

\vspace{2cm}
\hfill ———myp
